\documentclass[12pt,a4paper]{article}
\usepackage[utf8]{inputenc}
\usepackage[margin=1in]{geometry}
\usepackage{graphicx}
\usepackage{setspace}
\usepackage{parskip}
\usepackage{booktabs}
\usepackage{enumitem}
\usepackage{hyperref}

\hypersetup{
    colorlinks=true,
    linkcolor=black,
    citecolor=black,
    urlcolor=blue
}

\setstretch{1.2}

\begin{document}

% --- TITLE PAGE ---
\begin{titlepage}
    \centering
    \setstretch{1.3}
    {\fontsize{20}{24}\selectfont \textbf{BANARAS HINDU UNIVERSITY}}\\[0.3cm]
    {\fontsize{18}{22}\selectfont \textbf{INSTITUTE OF SCIENCE}}\\[0.2cm]
    {\fontsize{14}{18}\selectfont \textbf{Department of Mathematics}}\\[1cm]
    
    \begin{figure}[h!]
        \centering
        \begin{minipage}{0.38\textwidth}
            \centering
            \includegraphics[height=3.5cm]{science_logo.png}
        \end{minipage}%
        \begin{minipage}{0.38\textwidth}
            \centering
            \includegraphics[height=3.5cm]{bhu_logo.png}
        \end{minipage}
    \end{figure}
    
    \vspace{1.2cm}
    
    {\fontsize{22}{26}\selectfont \textbf{Mathematical Renaissance}}\\[0.4cm]
    {\fontsize{15}{18}\selectfont \textit{(A Model for India)}}\\[1.5cm]
    
    \begin{minipage}{0.45\textwidth}
        \setstretch{1.2}
        \textit{Submitted by:}\\
        \textbf{Akshit Agarwal}\\
        Roll No: 24220MAT024\\
        Enrollment No: 475073
    \end{minipage}%
    \begin{minipage}{0.45\textwidth}
        \setstretch{1.2}
        \textit{Guided by:}\\
        \textbf{Prof. Ashish Kumar}\\
        \textbf{Upadhyay Sir}
    \end{minipage}
    
    \vfill
    
    {\fontsize{12}{14}\selectfont Academic Year 2025-2026}\\[0.2cm]
    {\fontsize{12}{14}\selectfont May 15, 2026}
\end{titlepage}

\newpage
\pagenumbering{roman}

% --- ACKNOWLEDGEMENT ---
\addcontentsline{toc}{section}{Acknowledgement}
\begin{center}
    {\Large \textbf{Acknowledgement}}\\[0.3cm]
    \textit{May 15, 2026}
\end{center}

\vspace{0.8cm}

I would like to express my sincere gratitude to my respected supervisor and professor, \textbf{Dr. Ashish Kumar Upadhyay}, for his valuable guidance and encouragement throughout the completion of this internship report on \textit{“A Mathematical model for India’s Mathematical Renaissance”}. His suggestions and motivation greatly helped me in understanding the topic and completing the work successfully.

\newpage

% --- TABLE OF CONTENTS ---
\cleardoublepage
\addcontentsline{toc}{section}{Table of Contents}
\tableofcontents
\newpage

\pagenumbering{arabic}

% --- SECTION 1: INTRODUCTION ---
\section{Introduction}

Biologically humans are born with only two fears, the fear of loud noises and fear of falling, rest other fears are a response to external stimuli and taught by people around us. The fear of reptiles, arachnids and others are somewhat induced in us.

One such fear among people is the fear of Mathematics. The fear of mathematics, when logically analyzed is merely irrational and arise due to misunderstanding and poor knowledge of intricacies and beauty of Mathematics.

In the article, \textit{“The nation that fears Numbers”}, the author provided qualitative diagnosis of India’s Civilizational phobia of Mathematics. This could also be traced back to the later Vedic period, when females and lower classes of society were not allowed the privilege of education thus developing lower scientific attitude among people, which further (a possibility) may have lead to the development of fear of mathematics. As suggested by a scientific study, a newborn inherits their thinking potential from their mother, which explains a lot. So, the anxiety generated from mathematics is not developed, but passed on from generations to generations as a non-linear dynamical system. When a student asks their mentor for a query but are left with a not that satisfactory explanation, a psychological statement could be stuck in their head: \textit{“Is Mathematics really that tough”}.

This analysis provides us with the fact that phobia of mathematics is not internally developed but is transmitted from our mentors like teachers, parents, classmates, peers, seniors and others. The system of mathematical anxiety has two stable equilibria: a \textbf{low-mathematics trap} (the current state, in which anxiety self-perpetuates through teacher transmission and parental messaging) and a \textbf{high-mathematics attractor} (the desired state, in which mathematical confidence self-reinforces). We characterize the unstable saddle point (the tipping threshold) separating the two basins of teacher transformation, attraction, and we identify the three intervention classes that can shift the system early childhood investment, and institutional revaluation across the tipping threshold most efficiently. We derive quantitative targets, timelines, and investment requirements for each intervention, and construct a composite \textit{“Mathematical Health Index”} (MHI) for India that can be tracked annually and compared internationally. The model predicts that a sustained, correctly sequenced investment program can move India across the tipping threshold in 12--15 years and achieve a high-mathematics equilibrium within 25 years — the window required for India to lead rather than follow the mathematics-intensive technologies of the mid-21st century.

\newpage

% --- SECTION 2: THE QUANTITATIVE ACTION PLAN ---
\section{The Quantitative Action Plan}

\subsection{10 Concrete Actions in 10 Years}

\subsubsection{ASER + Anxiety Survey}
\begin{itemize}[leftmargin=*]
    \item \textbf{What is proposed?} Add a mathematics anxiety test to the annual ASER education survey. Measure: fear of maths, confidence, learning level, district-wise performance.
    \item \textbf{Why?} Many students fail in maths because of fear and anxiety, not intelligence. The government needs data to track improvement.
    \item \textbf{Goal:} Create a baseline measurement for all districts. Use it like a feedback system to monitor reforms.
    \item \textbf{Cost:} INR 50 crore (Year 1)
\end{itemize}

\subsubsection{500 Mathematics Transformation Centres}
\begin{itemize}[leftmargin=*]
    \item \textbf{What is proposed?} Create 500 teacher-training mathematics centres across India. Each centre trains about 200 teachers every year.
    \item \textbf{Training focus:} Conceptual understanding, Problem solving, Deep mathematical thinking, Residential 6-week programmes.
    \item \textbf{Why?} Poor maths teaching is one of the biggest causes of weak learning. Better teachers $\rightarrow$ better students.
    \item \textbf{Managed by:} University mathematics departments. Funded by DST-SERB.
    \item \textbf{Cost:} INR 2,500 crore capital, INR 300 crore/year operational.
\end{itemize}

\subsubsection{NCERT Curriculum Rewrite}
\begin{itemize}[leftmargin=*]
    \item \textbf{What is proposed?} Rewrite Classes 1–5 maths curriculum. Done by mathematicians, not only administrators.
    \item \textbf{Key change:} Reduce unnecessary syllabus by 30\%. Increase conceptual depth.
    \item \textbf{Why?} Current system emphasizes memorization. Students don’t understand concepts properly.
    \item \textbf{Goal:} Fewer topics but deeper understanding.
    \item \textbf{Cost:} INR 25 crore
\end{itemize}

\subsubsection{Diagnostic Assessment System}
\begin{itemize}[leftmargin=*]
    \item \textbf{What is proposed?} Create annual diagnostic tests for every primary student. Tablet-based + paper version.
    \item \textbf{Purpose:} Identify exact conceptual weaknesses.
    \item \textbf{Important idea:} The report says: \textit{“Output is not a grade; it is a repair prescription.”} Meaning: Test should help fix weaknesses, not just give marks.
    \item \textbf{Why important?} Traditional exams only rank students. This system would diagnose learning gaps early.
    \item \textbf{Cost:} INR 100 crore development, INR 30 crore/year operation.
\end{itemize}

\subsubsection{Triple DST Mathematics Funding}
\begin{itemize}[leftmargin=*]
    \item \textbf{What is proposed?} Increase mathematics research funding from: INR 120 crore $\rightarrow$ INR 400 crore annually. Fill 4,000 vacant faculty positions.
    \item \textbf{Why?} India underinvests in pure mathematics research. Lack of professors and researchers weakens innovation.
    \item \textbf{Goal:} Strengthen university mathematics ecosystem.
    \item \textbf{Cost:} INR 280 crore additional/year, INR 500 crore/year for faculty salaries.
\end{itemize}

\subsubsection{“I Am a Maths Person” National Campaign}
\begin{itemize}[leftmargin=*]
    \item \textbf{What is proposed?} National awareness campaign to change public attitude toward mathematics.
    \item \textbf{Main idea:} The report says maths anxiety is often socially created: parents say “maths is hard,” teachers spread fear, society labels children early.
    \item \textbf{Campaign target groups:} Parents, Teachers, and Students.
    \item \textbf{Inspired by:} “Beti Bachao Beti Padhao”
    \item \textbf{Goal:} Reduce transmission of maths fear by parents by 30\%.
    \item \textbf{Cost:} INR 200 crore
\end{itemize}

\subsubsection{Stream Destreaming}
\begin{itemize}[leftmargin=*]
    \item \textbf{What is proposed?} Make quantitative reasoning compulsory for all Class 11–12 students regardless of stream. Includes: logical thinking, statistics, quantitative reasoning.
    \item \textbf{Why?} Currently: Arts students avoid maths, Commerce students avoid advanced reasoning, Science students dominate quantitative fields. This creates division and fear.
    \item \textbf{Goal:} Basic mathematical literacy for everyone.
    \item \textbf{Cost:} Negligible
\end{itemize}

\subsubsection{Five New Mathematics Research Clusters}
\begin{itemize}[leftmargin=*]
    \item \textbf{What is proposed?} Create 5 major regional mathematics research centres.
    \item \textbf{Model:} Based on ICTS Bangalore.
    \item \textbf{Functions:} Fundamental research, Teacher engagement, Collaboration, Advanced mathematics training.
    \item \textbf{Why?} India lacks world-class mathematics hubs.
    \item \textbf{Goal:} Create regional centres of excellence.
    \item \textbf{Cost:} INR 1,000 crore capital, INR 150 crore/year operation.
\end{itemize}

\subsubsection{Mathematician in Society Fellowship}
\begin{itemize}[leftmargin=*]
    \item \textbf{What is proposed?} 1,000 fellowships for mathematicians to work outside academia.
    \item \textbf{Areas:} Government ministries, Media organizations, Industry.
    \item \textbf{Roles:} Data analysis, Quantitative policy work, Mathematical consulting.
    \item \textbf{Why?} Mathematics should influence real-world decision-making.
    \item \textbf{Goal:} Make mathematical thinking visible in society.
    \item \textbf{Cost:} INR 200 crore/year
\end{itemize}

\subsubsection{Fields Medal Preparatory Track}
\begin{itemize}[leftmargin=*]
    \item \textbf{What is proposed?} Identify 50 exceptionally talented students every year after Class 11.
    \item \textbf{Provide:} Olympiad pathway, Elite mentoring and PhD/postdoctoral guidance.
    \item \textbf{Goal:} Prepare India for its first Fields Medal within 20 years.
    \item \textbf{What is Fields Medal?} The highest international award in mathematics (similar to Nobel Prize for maths).
    \item \textbf{Cost:} INR 20 crore/year
\end{itemize}

\subsection{Financial Feasibility \& Budget Comparison}
India’s annual education budget = INR 1.48 lakh crore. Entire 10-year maths reform plan costs only: 1.4\% of one year’s education budget. Meaning: The author argues this transformation is financially very feasible.

\subsection{Total Investment and Return}
\begin{table}[h!]
    \centering
    \caption{Summary of Proposed Investments and Annual Operational Costs}
    \vspace{0.2cm}
    \begin{tabular}{lcc}
        \toprule
        \textbf{Program} & \textbf{Capital (Rs. Cr.)} & \textbf{Annual Opex (Rs. Cr.)} \\
        \midrule
        Teacher Transformation Centers & 2500 & 300 \\
        NCERT Curriculum Rewrite & 25 & -- \\
        Diagnostic Assessment System & 100 & 30 \\
        DST Mathematics Funding Increase & -- & 280 \\
        Faculty Hiring (4000 posts) & -- & 500 \\
        National Campaign & -- & 200 \\
        Mathematics Research Cluster (5) & 1000 & 150 \\
        Mathematician in Society Fellowship & -- & 200 \\
        Fields Medal Preparatory Track & -- & 20 \\
        ASER Anxiety Survey & 50 & 15 \\
        \midrule
        \textbf{TOTAL} & \textbf{3675} & \textbf{1695} \\
        \bottomrule
    \end{tabular}
\end{table}

\subsection{Overall Core Idea of the Document}
The author believes India’s mathematics crisis is NOT because students are incapable. Main problems are: maths anxiety, weak conceptual teaching, rote learning, poor teacher training, low research funding, social fear of mathematics. The plan tries to solve all of these together:
\begin{center}
\textit{Better teachers $\rightarrow$ Better curriculum $\rightarrow$ Better diagnostics $\rightarrow$ Better research ecosystem $\rightarrow$ Better public attitude toward maths}
\end{center}

\subsection{Summary}
The thesis proposes a national mathematical renaissance where India treats mathematics as a strategic national capability, similar to science, technology, or infrastructure.

\newpage

% --- SECTION 3: REFERENCES ---
\section{References}
\textit{May 15, 2026}

The content of this work was taken from the research paper \textit{“A Mathematical Model for India’s Mathematical Renaissance”}.

\end{document}
